% Example presentation using CSH theme
\documentclass[aspectratio=169]{beamer}
\usetheme{csh}

\title{Example Presentation}
\subtitle{Using the CSH Beamer Theme}
\author{Your Name}
\institute{Complexity Science Hub}
\date{\today}

\begin{document}

\begin{frame}
    \titlepage
\end{frame}

\begin{frame}{Overview}
    \tableofcontents
\end{frame}

\section{Introduction}

\begin{frame}{Welcome}
    This is an example presentation using the \textbf{CSH theme}.

    \begin{itemize}
        \item First point
        \item Second point
        \item Third point
    \end{itemize}
\end{frame}

\begin{frame}{Features}
    \begin{block}{Block Title}
        This is a block with some content.
    \end{block}

    \begin{alertblock}{Alert Block}
        This is an alert block.
    \end{alertblock}

    \begin{exampleblock}{Example Block}
        This is an example block.
    \end{exampleblock}
\end{frame}

\section{Content}

\begin{frame}{Enumeration and Math}
    \begin{enumerate}
        \item First item
        \item Second item
        \item Third item
    \end{enumerate}

    \vspace{1em}

    Some inline math: $E = mc^2$

    \[
        \sum_{i=1}^{n} i = \frac{n(n+1)}{2}
    \]
\end{frame}

\begin{frame}{Figures}
    \begin{figure}
        \centering
        \includegraphics[width=0.3\textwidth]{csh-logo-web}
        \caption{Example figure with caption}
    \end{figure}
\end{frame}

\begin{frame}{Tables}
    \begin{table}
        \centering
        \begin{tabular}{lcc}
            \hline
            \textbf{Method} & \textbf{Accuracy} & \textbf{Runtime} \\
            \hline
            Baseline & 0.72 & 1.2s \\
            Our approach & \textbf{0.89} & 0.8s \\
            \hline
        \end{tabular}
        \caption{Comparison of methods}
    \end{table}
\end{frame}

\begin{frame}{Two-Column Layout}
    \begin{columns}
        \begin{column}{0.48\textwidth}
            \textbf{Pros}
            \begin{itemize}
                \item Fast convergence
                \item Low memory usage
                \item Easy to implement
            \end{itemize}
        \end{column}
        \begin{column}{0.48\textwidth}
            \textbf{Cons}
            \begin{itemize}
                \item Limited scalability
                \item Requires tuning
                \item Sensitive to noise
            \end{itemize}
        \end{column}
    \end{columns}
\end{frame}

\begin{frame}[fragile]{Code Listings}
    \begin{verbatim}
def hello_world():
    print("Hello, World!")
    return 42
    \end{verbatim}
\end{frame}

\begin{frame}{Source Citation}
    You can add source citations (e.g., links to websites) to your slides.

    \source{Source: Your citation here}
\end{frame}

\begin{frame}{References}
    \begin{thebibliography}{9}
        \bibitem{thurner2018}
        S.~Thurner, R.~Hanel, and P.~Klimek.
        \textit{Introduction to the Theory of Complex Systems}.
        Oxford University Press, 2018.

        \bibitem{battiston2016}
        S.~Battiston et al.
        Complexity theory and financial regulation.
        \textit{Science}, 351(6275):818--819, 2016.
    \end{thebibliography}
\end{frame}

% Closing slide
\cshclosingslide{Your Name}{lastname@csh.ac.at}{https://csh.ac.at}

\end{document}
